% !TEX root = ../dabirestan.tex
%% ============================================================
%%  فصل ۴ — احتمال شرطی و نمودار درختی
%% ============================================================
\chapter{احتمال شرطی و نمودار درختی}
\label{hs:ch4}

\chapterintro{\begin{itemize}
  \item بفهمید «اطلاعات جدید» چطور احتمال را تغییر می‌دهد؛
  \item احتمال شرطی را با فرمول و با «کوچک کردن فضای نمونه» حساب کنید؛
  \item با قانون ضرب و نمودار درختی، احتمال آزمایش‌های چندمرحله‌ای را بیابید.
\end{itemize}}

%% ------------------------------------------------------------
\section{وقتی چیزی می‌دانیم}
%% ------------------------------------------------------------

\begin{example}[label=hs:ex:hint]{یک راهنمایی کوچک}
دوستتان تاسی می‌اندازد و نتیجه را پنهان می‌کند. احتمال اینکه ۶ آمده باشد $\tfrac16$ است. حالا به شما می‌گوید: «عدد زوج آمده است.» حالا احتمال ۶ چقدر است؟
\begin{solution}
با این راهنمایی، برآمدهای فرد کنار می‌روند و فقط $\{2,4,6\}$ باقی می‌ماند. این سه حالت هنوز هم‌شانس‌اند و یکی از آن‌ها ۶ است؛ پس احتمال حالا $\tfrac13$ است. \emph{اطلاعات جدید، احتمال را عوض کرد.}
\end{solution}
\end{example}

این احتمالِ تازه را «احتمال ۶ \emph{به شرط} زوج بودن» می‌نامیم و می‌نویسیم $\P(\fa{۶}\mid\fa{زوج})$. خط عمودی $\mid$ را «به شرط» می‌خوانیم.

\begin{figblock}
\begin{tikzpicture}[scale=0.85]
  \draw[graytext, thick, rounded corners=5pt] (-2,-1.5) rectangle (3.2,1.5);
  \fill[setA, opacity=0.25] (0,0) circle (1.1); \fill[setB, opacity=0.3] (1.3,0) circle (1.1);
  \draw[setA!80!black, thick] (0,0) circle (1.1); \draw[setB!80!black, thick] (1.3,0) circle (1.1);
  \node at (-0.45,0.2) {\lr{$A$}}; \node at (1.85,0.2) {\lr{$B$}}; \node at (2.9,1.2) {\lr{$S$}};
  \node[font=\small] at (0.6,-1.95) {\rl{قبل از اطلاع: کل $S$}};
  \draw[-{Stealth[length=3mm]}, line width=1.1pt, accent] (3.6,0) -- node[above, font=\small]{\rl{$B$ رخ داد}} (5.4,0);
  \begin{scope}[xshift=7.4cm]
    \draw[graytext!40, thick, dashed, rounded corners=5pt] (-2,-1.5) rectangle (3.2,1.5);
    \begin{scope}\clip (1.3,0) circle (1.1); \fill[setA!70] (0,0) circle (1.1); \end{scope}
    \draw[setA!40, thick, dashed] (0,0) circle (1.1);
    \draw[setB!80!black, line width=1.6pt] (1.3,0) circle (1.1);
    \node[font=\scriptsize] at (0.62,0) {\lr{$A\!\cap\!B$}}; \node at (2.05,0.35) {\lr{$B$}};
    \node[font=\small] at (0.6,-1.95) {\rl{بعد از اطلاع: $B$ فضای نمونهٔ جدید است}};
  \end{scope}
\end{tikzpicture}

\figcaption[hs:fig:cond]{احتمال شرطی یعنی فضای نمونه را به $B$ کوچک کنیم و ببینیم چه سهمی از $B$ در $A$ هم هست.}
\end{figblock}

%% ------------------------------------------------------------
\section{تعریف احتمال شرطی}
%% ------------------------------------------------------------

\begin{definition}[label=hs:def:cond]{احتمال شرطی}
اگر $\P(B)>0$، احتمال $A$ به شرط $B$ برابر است با
\[ \P(A\mid B)=\frac{\P(A\cap B)}{\P(B)}. \]
اگر برآمدها هم‌شانس باشند، این همان است که بگوییم \ $\P(A\mid B)=\dfrac{n(A\cap B)}{n(B)}$.
\end{definition}

\begin{example}[label=hs:ex:glasses]{جدول دانش‌آموزان}
در کلاسی ۳۰ نفره وضعیت عینکی بودن دانش‌آموزان این است:
\begin{center}
\renewcommand{\arraystretch}{1.35}
\begin{tabular}{l|cc|c}
 & عینکی & بدون عینک & جمع\\\hline
دختر & $6$ & $10$ & $16$\\
پسر & $4$ & $10$ & $14$\\\hline
جمع & $10$ & $20$ & $30$
\end{tabular}
\end{center}
یک نفر به تصادف انتخاب می‌شود. احتمال‌های زیر را بیابید: (الف) عینکی باشد؛ (ب) عینکی باشد به شرط اینکه دختر باشد؛ (ج) دختر باشد به شرط اینکه عینکی باشد.
\begin{solution}
(الف) $\P(\fa{عینکی})=\tfrac{10}{30}=\tfrac13$.\\
(ب) فقط ۱۶ دختر را در نظر می‌گیریم که ۶ نفرشان عینکی‌اند: $\P(\fa{عینکی}\mid\fa{دختر})=\tfrac{6}{16}=\tfrac38$.\\
(ج) فقط ۱۰ عینکی را در نظر می‌گیریم که ۶ نفرشان دخترند: $\P(\fa{دختر}\mid\fa{عینکی})=\tfrac{6}{10}=\tfrac35$.
\end{solution}
\end{example}

\begin{warning}
$\P(A\mid B)$ با $\P(B\mid A)$ فرق دارد! در مثال بالا $\P(\fa{عینکی}\mid\fa{دختر})=\tfrac38$ ولی $\P(\fa{دختر}\mid\fa{عینکی})=\tfrac35$. همیشه دقت کنید \emph{کدام پیشامد شرط است} (بعد از خط عمودی).
\end{warning}

\begin{example}[label=hs:ex:dice-cond]{دو تاس با شرط}
دو تاس سالم پرتاب می‌شوند و می‌دانیم مجموع آن‌ها ۸ شده است. احتمال اینکه یکی از تاس‌ها ۶ باشد چقدر است؟
\begin{solution}
برآمدهای با مجموع ۸: $\{(2,6),(3,5),(4,4),(5,3),(6,2)\}$؛ یعنی $n(B)=5$. از این‌ها دو تا ۶ دارند. پس احتمال $\tfrac25$ است. (بدون اطلاع، احتمال اینکه دست‌کم یک تاس ۶ باشد $\tfrac{11}{36}\approx0.31$ بود؛ حالا $0.4$ شد.)
\end{solution}
\end{example}

%% ------------------------------------------------------------
\section{قانون ضرب}
%% ------------------------------------------------------------

اگر فرمول احتمال شرطی را در $\P(B)$ ضرب کنیم، قانونی به دست می‌آید که برای آزمایش‌های چندمرحله‌ای بسیار مفید است:

\begin{theorem}[label=hs:thm:mult]{قانون ضرب}
\[ \P(A\cap B)=\P(A)\times\P(B\mid A). \]
یعنی: احتمال اینکه «اول $A$ و بعد $B$» رخ دهد، برابر است با احتمال $A$ ضرب در احتمال $B$ وقتی می‌دانیم $A$ رخ داده.
\end{theorem}

\begin{example}[label=hs:ex:two-reds]{دو مهره پشت سر هم}
جعبه‌ای ۵ مهرهٔ قرمز و ۳ مهرهٔ آبی دارد. دو مهره را \emph{یکی پس از دیگری} و \emph{بدون برگرداندن} برمی‌داریم. احتمال اینکه هر دو قرمز باشند چقدر است؟
\begin{solution}
احتمال قرمز بودن اولی $\tfrac58$ است. اگر اولی قرمز باشد، ۴ قرمز از ۷ مهره باقی می‌ماند؛ پس احتمال قرمز بودن دومی به شرط قرمز بودن اولی $\tfrac47$ است. طبق قانون ضرب:
\[ \P(\fa{هر دو قرمز})=\frac58\times\frac47=\frac{20}{56}=\frac5{14}. \]
(این همان جواب مثال \ref{hs:ex:urn} در فصل ۳ است که با ترکیب حساب کردیم.)
\end{solution}
\end{example}

\begin{example}[label=hs:ex:aces]{دو آس}
از یک دست ۵۲ کارتی، دو کارت پشت سر هم (بدون برگرداندن) برمی‌داریم. احتمال اینکه هر دو آس باشند چقدر است؟
\begin{solution}
$\P=\dfrac{4}{52}\times\dfrac{3}{51}=\dfrac{12}{2652}=\dfrac{1}{221}$.
\end{solution}
\end{example}

%% ------------------------------------------------------------
\section{نمودار درختی}
%% ------------------------------------------------------------

نمودار درختی بهترین ابزار برای دیدن همهٔ حالت‌های یک آزمایش چندمرحله‌ای است. روی هر شاخه احتمالِ آن مرحله (به شرط مراحل قبل) را می‌نویسیم.

\begin{figblock}
\begin{tikzpicture}[
  grow=left, level distance=3.1cm,
  level 1/.style={sibling distance=2.5cm},
  level 2/.style={sibling distance=1.15cm},
  every node/.style={font=\small},
  root/.style={circle, fill=primary, inner sep=2pt},
  st/.style={draw=primary!60, fill=primarylight, rounded corners=2pt, inner sep=3pt},
  edge from parent/.style={draw=primary!55, thick},
  lab/.style={font=\scriptsize, fill=white, inner sep=1pt}]
  \node[root] {}
    child { node[st] {\rl{قرمز}}
      child { node[st] (rr) {\rl{قرمز}} edge from parent node[lab, above] {\lr{$\frac{4}{7}$}} }
      child { node[st] (rb) {\rl{آبی}} edge from parent node[lab, below] {\lr{$\frac{3}{7}$}} }
      edge from parent node[lab, above] {\lr{$\frac{5}{8}$}} }
    child { node[st] {\rl{آبی}}
      child { node[st] (br) {\rl{قرمز}} edge from parent node[lab, above] {\lr{$\frac{5}{7}$}} }
      child { node[st] (bb) {\rl{آبی}} edge from parent node[lab, below] {\lr{$\frac{2}{7}$}} }
      edge from parent node[lab, below] {\lr{$\frac{3}{8}$}} };
  \node[anchor=east, font=\small, text=accent!80!black] at ([xshift=-4mm]rr.west) {\lr{$\frac58\cdot\frac47=\frac{20}{56}$}};
  \node[anchor=east, font=\small] at ([xshift=-4mm]rb.west) {\lr{$\frac58\cdot\frac37=\frac{15}{56}$}};
  \node[anchor=east, font=\small] at ([xshift=-4mm]br.west) {\lr{$\frac38\cdot\frac57=\frac{15}{56}$}};
  \node[anchor=east, font=\small] at ([xshift=-4mm]bb.west) {\lr{$\frac38\cdot\frac27=\frac{6}{56}$}};
  \node[font=\scriptsize, text=graytext] at (-3.1,2.1) {\rl{مهرهٔ اول}};
  \node[font=\scriptsize, text=graytext] at (-6.2,2.1) {\rl{مهرهٔ دوم}};
\end{tikzpicture}

\figcaption[hs:fig:tree-urn]{نمودار درختی برداشتن دو مهره از جعبه‌ای با ۵ قرمز و ۳ آبی (بدون برگرداندن).}
\end{figblock}

\begin{theorem}[label=hs:thm:tree]{قاعده‌های نمودار درختی}
\begin{enumerate}
  \item احتمال شاخه‌هایی که از یک نقطه جدا می‌شوند، جمعشان ۱ است.
  \item احتمال هر مسیر = \textbf{ضرب} احتمال‌های روی شاخه‌های آن (قانون ضرب).
  \item احتمال یک پیشامد = \textbf{جمع} احتمال همهٔ مسیرهایی که به آن پیشامد می‌رسند.
\end{enumerate}
\end{theorem}

\begin{example}[label=hs:ex:same-color]{هم‌رنگ}
در شکل \ref{hs:fig:tree-urn}، احتمال اینکه دو مهره هم‌رنگ باشند چقدر است؟ احتمال اینکه مهرهٔ دوم قرمز باشد چطور؟
\begin{solution}
هم‌رنگ: دو مسیر «قرمز-قرمز» و «آبی-آبی»: \ $\tfrac{20}{56}+\tfrac{6}{56}=\tfrac{26}{56}=\tfrac{13}{28}$.\\
مهرهٔ دوم قرمز: دو مسیر «قرمز-قرمز» و «آبی-قرمز»: \ $\tfrac{20}{56}+\tfrac{15}{56}=\tfrac{35}{56}=\tfrac58$. جالب است که همان احتمال قرمز بودن مهرهٔ اول است! (اگر ندانیم اولی چه بوده، دومی هیچ فرقی با اولی ندارد.)
\end{solution}
\end{example}

\begin{example}[label=hs:ex:exam]{قبولی در دو مرحله}
در یک مسابقه، ۴۰٪ شرکت‌کنندگان از مرحلهٔ اول قبول می‌شوند و از قبول‌شدگان مرحلهٔ اول، ۲۵٪ در مرحلهٔ دوم هم قبول می‌شوند. احتمال اینکه یک شرکت‌کننده هر دو مرحله را قبول شود چقدر است؟
\begin{solution}
$\P(\fa{قبول اول و دوم})=\P(\fa{قبول اول})\times\P(\fa{قبول دوم}\mid\fa{قبول اول})=0.4\times0.25=0.1$؛ یعنی ۱۰٪.
\end{solution}
\end{example}

%% ------------------------------------------------------------
\section*{خلاصهٔ فصل ۴}
\addcontentsline{toc}{section}{خلاصه و تمرین‌ها}
%% ------------------------------------------------------------

\begin{summary}[title={\leavevmode\boxicon{accent}{$\Sigma$}\ \ خلاصهٔ فصل}]
\begin{itemize}
  \item \textbf{احتمال شرطی:} $\P(A\mid B)=\frac{\P(A\cap B)}{\P(B)}$؛ یعنی $B$ فضای نمونهٔ جدید می‌شود. در حالت هم‌شانس: $\frac{n(A\cap B)}{n(B)}$.
  \item $\P(A\mid B)$ و $\P(B\mid A)$ معمولاً برابر نیستند.
  \item \textbf{قانون ضرب:} $\P(A\cap B)=\P(A)\,\P(B\mid A)$.
  \item \textbf{نمودار درختی:} روی شاخه‌ها ضرب، بین مسیرها جمع.
\end{itemize}
\end{summary}

\subsection*{تمرین‌ها}

\begin{exercise}
تاسی سالم پرتاب می‌شود. احتمال اینکه عدد ۲ آمده باشد، به شرط اینکه بدانیم عدد کمتر از ۴ است، چقدر است؟ به شرط اینکه بدانیم عدد اول است چطور؟
\answer{$\tfrac13$؛ \ $\tfrac13$}
\end{exercise}

\begin{exercise}
در مثال \ref{hs:ex:glasses}، احتمال اینکه دانش‌آموز انتخاب‌شده پسر باشد به شرط اینکه عینکی نباشد چقدر است؟
\answer{$\tfrac{10}{20}=\tfrac12$}
\end{exercise}

\begin{exercise}
دو کارت پشت سر هم و بدون برگرداندن از دست ۵۲ کارتی برمی‌داریم. احتمال اینکه هر دو دل باشند چقدر است؟
\answer{$\tfrac{13}{52}\times\tfrac{12}{51}=\tfrac1{17}$}
\end{exercise}

\begin{exercise}
جعبه‌ای ۳ مهرهٔ سفید و ۲ مهرهٔ سیاه دارد. دو مهره پشت سر هم و بدون برگرداندن برمی‌داریم. با رسم نمودار درختی احتمال‌های (الف) هر دو سفید؛ (ب) یکی سفید و یکی سیاه را بیابید.
\answer{(الف) $\tfrac3{10}$ \quad (ب) $\tfrac35$}
\end{exercise}

\begin{exercise}
اگر $\P(A)=0.6$، $\P(B)=0.5$ و $\P(A\cap B)=0.3$، مقادیر $\P(A\mid B)$ و $\P(B\mid A)$ را بیابید.
\answer{$0.6$ و $0.5$}
\end{exercise}

\begin{exercise}
۷۰٪ دانش‌آموزان یک مدرسه در کلاس تقویتی شرکت می‌کنند. از شرکت‌کنندگان در کلاس تقویتی ۸۰٪ و از بقیه ۵۰٪ در امتحان نمرهٔ بالای ۱۵ می‌گیرند. احتمال اینکه دانش‌آموزی هم در کلاس شرکت کرده باشد و هم نمرهٔ بالای ۱۵ بگیرد چقدر است؟
\answer{$0.7\times0.8=0.56$}
\end{exercise}
